\documentclass[conference]{IEEEtran}

\usepackage{graphicx}
\usepackage{amsmath,amssymb}
\usepackage{float}
\usepackage{booktabs}
\usepackage{multirow}

\begin{document}

\title{
    \textbf{Benchmark Report: 531.deepsjeng\_r} \\
    \text{Register Spill Analysis on RISC-V via gem5}
}

\author{
    \IEEEauthorblockN{Lütfullah Burak Kaya}
    \IEEEauthorblockA{
        Ozyegin University \\
        burak.kaya.50711@ozu.edu.tr
    }
    \and
    \IEEEauthorblockN{Ismail Akturk}
    \IEEEauthorblockA{
        Ozyegin University \\
        ismail.akturk@ozyegin.edu.tr
    }
}

\newcommand{\LongDate}{%
    \number\day\space
    \ifcase\month
    \or January\or February\or March\or April\or May\or June%
    \or July\or August\or September\or October\or November\or December%
    \fi
    \space \number\year
}

\twocolumn[
\begin{@twocolumnfalse}
\maketitle
\begin{center}{\small \LongDate}\end{center}
\vspace{1em}
\end{@twocolumnfalse}
]

% ================================================

\begin{abstract}
This report presents the register spill characterization of the
\texttt{531.deepsjeng\_r} benchmark from SPEC CPU2017, executed on a
RISC-V (RV64GC) target using the gem5 TimingSimpleCPU simulator.
A custom SpillDetector was integrated into gem5 to count store-then-load
patterns on the stack within a user-defined Region of Interest (ROI).
The test dataset (a single chess position, depth\,15) was used as the
input. Results show a spill rate of \textbf{4.17\%} within the ROI,
corresponding to 1.93 billion detected spill events across 46.2 billion
ROI instructions. The alpha-beta search engine creates deep, recursive
call stacks that maintain large amounts of board state simultaneously,
driving persistent register pressure throughout the search.
\end{abstract}

% =========================================================
\section{Benchmark Overview}

\texttt{531.deepsjeng\_r} is the SPEC CPU2017 rate variant of Deep
Sjeng, a computer chess engine implementing minimax search with
alpha-beta pruning, quiescence search, and move ordering heuristics.
Given a chess position in Forsyth-Edwards Notation (FEN) and a search
depth, Deep Sjeng searches all legal continuations to the specified
depth, evaluating leaf positions with a static evaluation function.
The workload exercises recursive tree traversal with irregular control
flow, frequent function calls, transposition table lookups, and
move-ordering bookkeeping.

The RISC-V binary (\texttt{deepsjeng\_r.riscv}) was compiled with gem5
ROI markers inserted in \texttt{src/sjeng.cpp}, wrapping the
\texttt{run\_epd\_testsuite()} function call from \texttt{main()}:

\begin{itemize}
    \item \texttt{m5\_work\_begin(0,0)} --- before \texttt{run\_epd\_testsuite()}
    \item \texttt{m5\_work\_end(0,0)}   --- after  \texttt{run\_epd\_testsuite()}
\end{itemize}

\texttt{run\_epd\_testsuite()} reads the EPD file (one position per
line), parses each FEN record, and runs a full alpha-beta search for
each position. This function encompasses all benchmark computation.

\subsection{Input Datasets}

SPEC CPU2017 provides three dataset sizes for \texttt{deepsjeng\_r},
differing in the number of chess positions and the required search depth
(Table~\ref{tab:inputs}).

\begin{table}[H]
\centering
\caption{deepsjeng\_r Input Dataset Sizes}
\label{tab:inputs}
\begin{tabular}{lcc}
\toprule
\textbf{Dataset} & \textbf{Positions} & \textbf{Search depth} \\
\midrule
test    & 1  & 15 \\
train   & 1  & 15 \\
refrate & 24 & 15 \\
\bottomrule
\end{tabular}
\end{table}

The \textbf{test} dataset was selected for this run. The single test
position is given in FEN notation:
\begin{verbatim}
r2qk2r/ppp1b1pp/2n1p3/3pP1n1/
3P2b1/2PB1NN1/PP4PP/R1BQK2R w KQkq - 0 1
\end{verbatim}
This is an open tactical position with both sides having active pieces,
requiring a deep combinatorial search. The search depth of 15 half-moves
(plies) results in an extremely large search tree, dominating the
simulation time.

% =========================================================
\section{Simulation Setup}

\subsection{gem5 Configuration}

\begin{itemize}
    \item \textbf{CPU model:}   TimingSimpleCPU
    \item \textbf{ISA:}         RISC-V RV64GC
    \item \textbf{Caches:}      enabled (default L1 I/D + L2)
    \item \textbf{Memory:}      4\,GB simulated physical RAM
    \item \textbf{Spill mode:}  summary only (\texttt{spill\_verbose=False})
\end{itemize}

The full gem5 invocation was:

\begin{verbatim}
./build/RISCV/gem5.opt
  --outdir=benchmarks/deepsjeng/m5out_test
  configs/deprecated/example/se.py
  --cmd=benchmarks/deepsjeng/deepsjeng_r.riscv
  --options="benchmarks/deepsjeng/test.txt"
  --cpu-type=TimingSimpleCPU
  --caches
  --mem-size=4GB
\end{verbatim}

% =========================================================
\section{Simulation Results}

The simulation completed successfully. Deep Sjeng searched the test
position to depth 15 and printed the best move sequence at each
iterative deepening step.

\subsection{Pre-ROI Context}

Before the ROI, the binary executed engine initialisation: loading
the transposition table, initialising the move generator, and parsing
the input file. This pre-ROI phase is exceptionally small compared to
the ROI workload. Table~\ref{tab:preroi} shows the global counters at
ROI entry.

\begin{table}[H]
\centering
\caption{Global Counters at ROI Entry}
\label{tab:preroi}
\begin{tabular}{lr}
\toprule
\textbf{Metric} & \textbf{Value} \\
\midrule
Total instructions (pre-ROI) & 115,433,501 \\
Total spills detected (pre-ROI) & 4,962 \\
Pre-ROI spill rate & 0.004\% \\
\bottomrule
\end{tabular}
\end{table}

The pre-ROI spill rate of 0.004\% is by far the lowest observed across
all benchmarks in this study. Engine initialisation involves mainly
linear table setup with simple sequential access patterns and minimal
register pressure, in stark contrast to the recursive search phase.

\subsection{ROI Spill Statistics}

Table~\ref{tab:roi} presents the spill statistics collected within
the ROI.

\begin{table}[H]
\centering
\caption{ROI Spill Statistics --- 531.deepsjeng\_r (test, depth\,15)}
\label{tab:roi}
\begin{tabular}{lr}
\toprule
\textbf{Metric} & \textbf{Value} \\
\midrule
ROI instructions         & 46,202,686,170 \\
ROI stores               &  3,395,185,975 \\
ROI loads                &  7,989,081,366 \\
\midrule
\textbf{ROI spills}      & \textbf{1,928,903,656} \\
\textbf{Spill rate}      & \textbf{4.1749\%} \\
\bottomrule
\end{tabular}
\end{table}

\subsection{Timing}

\begin{table}[H]
\centering
\caption{Simulation Timing}
\label{tab:timing}
\begin{tabular}{lr}
\toprule
\textbf{Metric} & \textbf{Value} \\
\midrule
Simulated time           & 78.15\,s \\
Host wall time           & 24{,}937\,s ($\approx$6.9\,hours) \\
Total simulated instructions & 46,318,246,441 \\
\bottomrule
\end{tabular}
\end{table}

% =========================================================
\section{Analysis}

\subsection{Spill Rate Interpretation}

The measured spill rate of \textbf{4.17\%} means that approximately
1 in every 24 instructions inside the ROI is a register spill.
Table~\ref{tab:compare} places deepsjeng among all evaluated benchmarks.

\begin{table}[H]
\centering
\caption{Cross-Benchmark Spill Rate Comparison}
\label{tab:compare}
\begin{tabular}{lrr}
\toprule
\textbf{Benchmark} & \textbf{ROI Spills} & \textbf{Spill Rate} \\
\midrule
507.cactuBSSN\_r & 3,782,099,438 & 7.8756\% \\
500.perlbench\_r &     1,070,055 & 7.2838\% \\
526.blender\_r   &    61,845,869 & 6.2193\% \\
541.leela\_r     & 1,158,236,271 & 4.4438\% \\
519.lbm\_r       &   274,130,440 & 4.2130\% \\
531.deepsjeng\_r & 1,928,903,656 & \textbf{4.1749\%} \\
544.nab\_r       &   151,310,616 & 2.2612\% \\
538.imagick\_r   &     2,424,899 & 2.0601\% \\
505.mcf\_r       &   446,113,136 & 1.8263\% \\
508.namd\_r      &     327,718,955 & 1.0269\% \\
\bottomrule
\end{tabular}
\end{table}

deepsjeng places sixth among the ten evaluated benchmarks, slightly
below lbm (4.21\%) and just above nab (2.26\%). Its spill rate is
consistent with other tree-search workloads (\texttt{541.leela\_r}:
4.44\%), reflecting the shared structural characteristic of deep
recursive call frames.

\subsection{Store vs.\ Load Balance}

The load count (7.99B) is approximately $2.35\times$ higher than the
store count (3.40B). This asymmetry is characteristic of alpha-beta
search: each internal node must read the current position's piece
lists, hash table entry, and move list to score and order moves, while
writes occur only when updating alpha/beta bounds, recording the best
move, or writing a new transposition table entry. The dominant access
pattern is read-heavy position evaluation.

\subsection{Spill Fraction of Loads}

Within the ROI, the ratio of spills to total loads is:
\[
\frac{1{,}928{,}903{,}656}{7{,}989{,}081{,}366} \approx 24.1\%
\]
Nearly 1 in 4 load instructions is a spill reload. Alpha-beta search
requires each recursive call to save and restore the current position's
state (side to move, piece lists, castling rights, en passant square,
move history) alongside the search parameters (alpha, beta, depth,
best-move candidate), creating register pressure at every call frame
that the RISC-V 32-register file cannot fully absorb.

\subsection{Pre-ROI Spill Density}

The pre-ROI spill rate of 0.004\% ($4{,}962$ spills over
$115{,}433{,}501$ instructions) is the lowest pre-ROI rate observed
across all benchmarks evaluated. Deep Sjeng's engine initialisation
consists mainly of sequential table fills and simple struct
assignments with few live variables simultaneously, resulting in
negligible register pressure before the search begins.

\subsection{ROI Coverage}

The ROI accounts for $46{,}202{,}686{,}170$ out of $46{,}318{,}246{,}441$
total simulated instructions:
\[
\frac{46{,}202{,}686{,}170}{46{,}318{,}246{,}441} \approx 99.75\%
\]
This is the highest ROI coverage fraction among all benchmarks evaluated,
confirming that virtually the entire simulation is occupied by the
alpha-beta search inside \texttt{run\_epd\_testsuite()}. The 0.25\%
outside the ROI corresponds to the trivially simple engine initialisation
phase.

\subsection{Search Tree Depth and Register Pressure}

Alpha-beta search at depth 15 creates call stacks up to 15 frames deep
in the main search routine, with additional frames for quiescence search
at tactical positions. At each frame, the engine must keep alive the
board state, move iterator, alpha and beta bounds, the current move
being examined, and the best-move record. With RISC-V providing 32
general-purpose registers, the compiler spills the excess live variables
to the stack at each recursive call and reloads them on return,
producing the observed 24.1\% spill fraction of loads.

The sharp contrast between the negligible pre-ROI spill rate (0.004\%)
and the high ROI rate (4.17\%) is the most distinctive characteristic
of deepsjeng: the benchmark transitions abruptly from a simple
initialisation sequence to an extremely register-intensive recursive
kernel with no gradual ramp-up.

% =========================================================
\section{Conclusion}

The \texttt{531.deepsjeng\_r} benchmark exhibits a register spill rate
of 4.17\% within its alpha-beta search ROI on RISC-V RV64GC. With 1.93
billion spill events across 46.2 billion ROI instructions, deepsjeng
places sixth among the ten evaluated SPEC CPU2017 benchmarks. The
2.35:1 load-to-store ratio and 24.1\% spill fraction of loads both
reflect the read-dominated nature of minimax tree traversal, where
position state is frequently read for move generation and evaluation
but written only when bounds or best-moves change. The most distinctive
characteristic of deepsjeng is the extreme contrast between its
negligible pre-ROI spill rate (0.004\%) --- the lowest of all benchmarks
--- and its substantial ROI spill rate, demonstrating that the chess
engine transitions from a trivial initialisation phase directly into a
highly register-pressure-intensive recursive search kernel.

\end{document}
